\section{Addressing the Problems}
\label{sec:addressing}

\Cref{sec:problems} listed the problems. This section describes the mechanism
chosen for each one, the alternatives that were set aside, and what each choice
costs.

\subsection{Mutual exclusion on a connector (P1)}
\label{sec:p1}

Each physical connector is owned by one Erlang process, a \texttt{gen\_statem}
implemented in \texttt{vs\_connector}. Every request that touches that outlet
becomes a message in that process's mailbox, and messages are handled one at a
time, so two drivers reserving the same connector in the same millisecond are
serialised by the runtime and the second finds it already \texttt{held}.

\paragraph{Which driver wins.} Since Erlang guarantees FIFO order point to point
and nothing across senders, the connector goes to the driver whose message the
runtime delivers first. Mutual exclusion is the property claimed here, and nothing
in this domain asks for more.

The machine gives a second guarantee: with the wrong vehicle, \texttt{charging} is
unreachable from \texttt{held}, so no session can start on a connector somebody
else holds. The transition does not exist, so there is no check to forget.

\subsection{One reservation per vehicle, network-wide (P2)}
\label{sec:p2}

The mechanism adopted is a
\emph{claim}: a permission, granted by the coordinator, recording that a given
vehicle is now committed and to which station.

\subsubsection{Claim first, commit after}

A station obtains the claim \emph{before} it creates the reservation.

In the wrong order, a crash between the two steps leaves a reservation that no
other node knows about, allowing the vehicle can then obtain a second reservation
elsewhere. In the right order, the same crash leaves a claim that was granted and
never used. Nobody can reserve for that vehicle for a few minutes, and then the
claim expires on its own. The system errs towards refusing rather than towards
granting twice.

\subsubsection{Why a central arbiter}

The coordinator is a single process, \texttt{vs\_coord\_srv}, and its claim table
is a map keyed by vehicle. Requests from different stations queue in its mailbox
and are decided one at a time. This is the centralised mutual exclusion where there is one critical section per vehicle. Two claims for two different cars
contend over nothing, and what the mailbox serialises is the instant of the
decision.

Three designs were considered.

\paragraph{One coordinator, unreplicated.} The simplest, and it leaves an
avoidable single point of failure.

\paragraph{A \texttt{UNIQUE} constraint in MySQL.} Correct, and almost free to
write. But the coordinator has to exist anyway for the cluster map and for node
monitoring, so this removes no component, and a claim that \emph{expires by
itself} is a timer inside a process, while in a database it becomes a cleanup job
or a filter on every read.

\paragraph{Direct coordination between stations (Ricart-Agrawala).} No central
authority, at $2(n-1)$ messages per reservation, and it serialises access to
\emph{one} critical section: here two reservations for two different vehicles
would wait for each other with nothing to gain.

What we built is the first one with its single point of failure removed: one
arbiter serving at a time, elected among three coordinator nodes and required to
hold a quorum, which is \cref{sec:p2b}.

\subsubsection{Ordering decided by one clock}

When a claim is granted, the coordinator stamps it with \texttt{granted\_at} and
sends that timestamp back with the grant. The station stores it and echoes it in
the renewal it sends every ten seconds for all the claims it holds
(\cref{sec:p3}).

The purpose is conflict resolution after a failover. If two claims for the same
vehicle ever meet in the table, the older one wins, and the comparison is between
two timestamps produced by the \emph{same} process. No two machines with
unsynchronised clocks are ever compared.

\subsection{Leases (P3)}
\label{sec:p3}

The problem of the driver who reserves and never arrives is approached by a lease: every reservation is granted with an expiry, and
when it passes the connector returns to \texttt{free} with no operator action and
no cooperation from the party that disappeared.

There are two nested expiries, and they protect against two different absent
parties:

\begin{itemize}
  \item the \textbf{reservation lease} protects against the driver who never
        shows up. It is the timer the connector process holds;
  \item the \textbf{claim expiry} protects against the station that dies without
        releasing what it holds. A renewal keeps the claim alive, so a station
        that crashes stops renewing and its claims expire on their own.
\end{itemize}


\subsection{Coordinator replication: election and quorum (P2b)}
\label{sec:p2b}

Three coordinator nodes run, \texttt{vs@coord1}, \texttt{vs@coord2} and
\texttt{vs@coord3}, and every node of the cluster is configured with the same
three names. One of them is the leader and is the only node that grants or refuses
claims. A station that addresses one of the other two is redirected to the leader.

\paragraph{Election.} When the leader stops responding, the survivors elect a new
one with the bully algorithm over those names: the list is sorted, and the highest
live entry wins, so a cluster in good health is led by \texttt{vs@coord3}. A node that notices the leader is gone sends \texttt{ELECTION}
to the nodes above it. Silence means it is the most senior node alive and it
announces \texttt{LEADER} downwards; an answer means somebody more senior is
handling it, and it waits. With three nodes this settles in two rounds of
messages. Detection is by timeout, which means we assume a \emph{synchronous}
system, an upper bound on response time.

\paragraph{Quorum.} Bully alone does not prevent split brain. In a partition each
side happily elects its own local maximum, and two authorities granting claims
break precisely the invariant the coordinator exists to protect. So a coordinator
serves only while it can see a majority of the configured coordinator cluster,
itself included: two out of three. A node in the minority suspends itself and
refuses everything, leader or not.

This is why the cluster has three members and not two. With an even number, a
partition can produce two halves that each consider themselves entitled to grant.

\paragraph{State after a failover.} Since a new leader inherits nothing, it asks the stations what they hold and
rebuilds the claim table from their answers. The reason of this approach is developed in
\cref{sec:rebuild}.

\subsection{Sharing the site power (P5)}
\label{sec:p5}

An Erlang module, \texttt{vs\_power}, recomputes the split on every arrival and departure, and on a
periodic tick that follows the charging curve as batteries fill.

The policy is max-min fair share with hand-back. Every active session starts from
$budget/N$. A session that cannot absorb its share, either because the car is
small or because it is tapering near full, takes only what it asks for and hands
the surplus back, which is then split again among the others. The loop repeats
until nobody hands anything back or the budget is exhausted.


\subsection{One source of truth for every client (P6)}
\label{sec:p6}

The cheapest way to make every client agree is to give them nothing to compute.
The station holds the authoritative state and pushes a \emph{complete} snapshot of
it, so that each page that misses a
message is wrong only until the next push, and one that reconnects is correct
immediately. Snapshots go out on every change and on a periodic tick.

